\documentclass[10pt,a4paper]{amsart}
\usepackage{amsmath,amssymb,amsthm,graphicx}
\usepackage{verbatim}
\newenvironment{reference}{\comment}{\endcomment}
%\newenvironment{reference}{}{}
\newenvironment{slogan}{\comment}{\endcomment}
\newenvironment{history}{\comment}{\endcomment}
% For commutative diagrams we use Xy-pic
\usepackage[all]{xy}

% We use 2cell for 2-commutative diagrams.
\xyoption{2cell}
% Theorem environments.
%
\theoremstyle{plain}
\newtheorem{theorem}[subsection]{Theorem}
\newtheorem{proposition}[subsection]{Proposition}
\newtheorem{lemma}[subsection]{Lemma}

\theoremstyle{definition}
\newtheorem{definition}[subsection]{Definition}
\newtheorem{example}[subsection]{Example}
\newtheorem{exercise}[subsection]{Exercise}
\newtheorem{situation}[subsection]{Situation}

\theoremstyle{remark}
\newtheorem{remark}[subsection]{Remark}
\newtheorem{remarks}[subsection]{Remarks}

\numberwithin{equation}{subsection}

% Macros
%
\def\lim{\mathop{\mathrm{lim}}\nolimits}
\def\colim{\mathop{\mathrm{colim}}\nolimits}
\def\Spec{\mathop{\mathrm{Spec}}}
\def\Hom{\mathop{\mathrm{Hom}}\nolimits}
\def\Ext{\mathop{\mathrm{Ext}}\nolimits}
\def\SheafHom{\mathop{\mathcal{H}\!\mathit{om}}\nolimits}
\def\SheafExt{\mathop{\mathcal{E}\!\mathit{xt}}\nolimits}
\def\Sch{\mathit{Sch}}
\def\Mor{\mathop{\mathrm{Mor}}\nolimits}
\def\Ob{\mathop{\mathrm{Ob}}\nolimits}
\def\Sh{\mathop{\mathit{Sh}}\nolimits}
\def\NL{\mathop{N\!L}\nolimits}
\def\CH{\mathop{\mathrm{CH}}\nolimits}
\def\proetale{{pro\text{-}\acute{e}tale}}
\def\etale{{\acute{e}tale}}
\def\QCoh{\mathit{QCoh}}
\def\Ker{\mathop{\mathrm{Ker}}}
\def\Im{\mathop{\mathrm{Im}}}
\def\Coker{\mathop{\mathrm{Coker}}}
\def\Coim{\mathop{\mathrm{Coim}}}

% Boxtimes
%
\DeclareMathSymbol{\boxtimes}{\mathbin}{AMSa}{"02}

%
% Macros for moduli stacks/spaces
%
\def\QCohstack{\mathcal{QC}\!\mathit{oh}}
\def\Cohstack{\mathcal{C}\!\mathit{oh}}
\def\Spacesstack{\mathcal{S}\!\mathit{paces}}
\def\Quotfunctor{\mathrm{Quot}}
\def\Hilbfunctor{\mathrm{Hilb}}
\def\Curvesstack{\mathcal{C}\!\mathit{urves}}
\def\Polarizedstack{\mathcal{P}\!\mathit{olarized}}
\def\Complexesstack{\mathcal{C}\!\mathit{omplexes}}
% \Pic is the operator that assigns to X its picard group, usage \Pic(X)
% \Picardstack_{X/B} denotes the Picard stack of X over B
% \Picardfunctor_{X/B} denotes the Picard functor of X over B
\def\Pic{\mathop{\mathrm{Pic}}\nolimits}
\def\Picardstack{\mathcal{P}\!\mathit{ic}}
\def\Picardfunctor{\mathrm{Pic}}
\def\Deformationcategory{\mathcal{D}\!\mathit{ef}}
\usepackage{cleveref}
\crefname{originalproposition}{original proposition}{original propositions}
\Crefname{originalproposition}{Original proposition}{Original propositions}
\makeatletter
\newlabel{etale-cohomology-proposition-finite-higher-direct-image-zero}{{03QP}{source}{}{sourcepointer-1}{}}
\newlabel{etale-cohomology-proposition-finite-higher-direct-image-zero@cref}{{[originalproposition][0][]03QP}{[1][0][]source}{}{}{}}
\makeatother
\crefname{originallemma}{original lemma}{original lemmas}
\Crefname{originallemma}{Original lemma}{Original lemmas}
\makeatletter
\newlabel{spaces-properties-lemma-pushforward-etale-base-change}{{03LR}{source}{}{sourcepointer-2}{}}
\newlabel{spaces-properties-lemma-pushforward-etale-base-change@cref}{{[originallemma][0][]03LR}{[1][0][]source}{}{}{}}
\makeatother
\title{The stalk of the pushforward along a finite morphism is the product of the stalks over the geometric lifts}
\author{The Stacks Project Authors}
\date{stacks-project:spaces-cohomology.tex:local-snapshot-20261012; GNU FDL 1.2 or later}
\begin{document}
\maketitle
\noindent\textit{Editorial note.} All delivered mathematical text is the author's own, copied from the pinned Stacks Project chapter file spaces-cohomology.tex; no statement or proof step was rewritten, reordered or removed.\par\medskip
\noindent\textit{Editorial note.} Preamble provenance: the pinned chapter file opens with \textbackslash{}input\{preamble\} and carries no preamble text of its own, so the project's single Stacks-wide preamble staged at sources/stacks-spaces-cohomology/preamble.tex is excerpted verbatim as three spans, lines 8-12 (verbatim package and the project's own reference/slogan/history display environments), lines 14-18 (Xy-pic setup) and lines 154-217 (theorem declarations and mathematical macros). The wrapper class and the amsmath/amssymb/amsthm/graphicx base packages are supplied by the assembler. No macro was invented or redefined.\par\medskip
\noindent\textit{Editorial note.} Same-chapter pointers are located at their \textbackslash{}label lines inside the pinned chapter file; cross-chapter pointers are located in the project's own tag list staged in the snapshot (sources/stacks-chow/tags.txt), at the verbatim lines carrying their canonical tags.\par\medskip
\noindent\textit{Allowed standard prerequisites (editorial).}\begin{itemize}
\item The stalk of a pushforward computed after an etale base change. Properties of Spaces, label spaces-properties-lemma-pushforward-etale-base-change (invoked at proof line 535). It lets the computation be performed after replacing \$Y\$ by an etale neighbourhood of the geometric point, and \$X\$ by the corresponding fibre product.
\item The statement for a finite morphism of schemes. Etale Cohomology, proposition etale-cohomology-proposition-finite-higher-direct-image-zero (invoked at proof line 538). After the reduction to schemes it identifies the stalk of the pushforward with the product of the stalks at the finitely many lifts.
\end{itemize}\medskip
\ifdefined\hypertarget\hypertarget{sourcepointer-1}{}\fi
\noindent\textit{Source reference.} Original proposition 03QP: sources/stacks-chow/tags.txt, source lines 4902--4902. Navigation pointer for \textbackslash{}ref\{etale-cohomology-proposition-finite-higher-direct-image-zero\} used inside the counted span (line 539); the argument that consumes it is retained in the span.\par\medskip
\ifdefined\hypertarget\hypertarget{sourcepointer-2}{}\fi
\noindent\textit{Source reference.} Original lemma 03LR: sources/stacks-chow/tags.txt, source lines 4723--4723. Navigation pointer for \textbackslash{}ref\{spaces-properties-lemma-pushforward-etale-base-change\} used inside the counted span (line 535); the argument that consumes it is retained in the span.\par\medskip
\begin{lemma}
\label{lemma-stalk-push-finite}
Let $S$ be a scheme. Let $f : X \to Y$ be a finite morphism of algebraic
spaces over $S$. Let $\overline{y}$ be a geometric point of $Y$ with
lifts $\overline{x}_1, \ldots, \overline{x}_n$ in $X$. Then
$$
(f_*\mathcal{F})_{\overline{y}} =
\prod\nolimits_{i = 1, \ldots, n}
\mathcal{F}_{\overline{x}_i}
$$
for any sheaf $\mathcal{F}$ on $X_\etale$.
\end{lemma}

\begin{proof}
Choose an \'etale neighbourhood $(V, \overline{v})$ of $\overline{y}$.
Then the stalk $(f_*\mathcal{F})_{\overline{y}}$
is the stalk of $f_*\mathcal{F}|_V$ at $\overline{v}$.
By Properties of Spaces,
Lemma \ref{spaces-properties-lemma-pushforward-etale-base-change}
we may replace $Y$ by $V$ and $X$ by $X \times_Y V$.
Then $Z \to X$ is a finite morphism of schemes and the result is
\'Etale Cohomology, Proposition
\ref{etale-cohomology-proposition-finite-higher-direct-image-zero}.
\end{proof}

\end{document}
